\Needspace{8\baselineskip}
\Article{283}
Aus diesen Grundsätzen ergibt sich die folgende Methode, um alle eigentlichen
Darstellungen der binären Form
\[
\varphi=ptt+2qtu+ruu
\]
mit der Determinante $D$ durch die ternäre Form $f$ mit der
Determinante $\Delta$ zu finden.

I.\ Man bestimme alle verschiedenen (\textit{d.h.}, nicht äquivalenten)
Werte des Ausdrucks $\surd\Delta(p,\allowbreak -q,\allowbreak r)(\mathrm{mod.}\ D)$. Für den Fall, dass $\varphi$ eine primitive Form
und $\Delta$ zu $D$ teilerfremd ist, wurde dieses Problem oben (Art.\ 233)
gelöst. Die übrigen Fälle lassen sich sehr leicht darauf zurückführen;
der Kürze halber können wir dies hier aber nicht ausführlich erklären.
Wir bemerken nur, dass, wenn $\Delta$ zu $D$ teilerfremd ist, der Ausdruck
$\Delta(p,-q,r)$ nur dann ein quadratischer Rest von $D$ sein kann, wenn
$\varphi$ eine primitive Form ist. Setzt man nämlich
\[
\Delta p=BB-DA',\qquad -\Delta q=BB'-DB'',\qquad \Delta r=B'B'-DA
\]
so erhält man
\[
(DB''-\Delta q)^2=(DA'+\Delta p)(DA+\Delta r)
\]
Durch Ausmultiplizieren und Einsetzen von $qq-pr$ für $D$ folgt daraus
\[
(qq-pr)(B''B''-AA')-\Delta(Ap+2B''q+A'r)+\Delta\Delta=0
\]
Man schließt daraus leicht, dass, wenn $p,q,r$ einen gemeinsamen Teiler
hätten, dieser auch $\Delta\Delta$ teilen würde; dann könnte aber $\Delta$ nicht
zu $D$ teilerfremd sein. Also können $p,q,r$ keinen gemeinsamen Teiler
haben; mit anderen Worten, $\varphi$ ist eine primitive Form.

II.\ Bezeichnen wir die Anzahl dieser Werte mit $m$, und nehmen wir
an, unter ihnen seien $n$ Werte, die sich selbst entgegengesetzt sind
(wobei wir $n=0$ setzen, wenn es keine solchen gibt). Dann bilden die
übrigen $m-n$ Werte offenbar Paare entgegengesetzter Werte (da wir
voraussetzen, dass alle Werte vollständig vorliegen). Man verwerfe aus
jedem Paar entgegengesetzter Werte einen beliebigen; insgesamt bleiben
dann $\frac{1}{2}(m+n)$ Werte übrig. So sind \textit{z.B.}\ von den acht Werten des Ausdrucks
$\surd-1(19,-3,41)$
\SourceLeaf{332}
$(\mathrm{mod.}\,770)$, nämlich $(39,237)$, $(171,-27)$, $(269,-83)$, $(291,-127)$, $(-39,-237)$, $(-171,27)$, $(-269,83)$,
$(-291,127)$, die letzten vier als den ersten vier entgegengesetzt zu verwerfen.
Ist ferner $(B,B')$ ein sich selbst entgegengesetzter Wert, so sind offenbar
$2B$, $2B'$ und daher auch $2\Delta p$, $2\Delta q$, $2\Delta r$ durch $D$
teilbar. Sind also $\Delta$, $D$ teilerfremd, so sind auch $2p$,
$2q$, $2r$ durch $D$ teilbar. Da nach I in diesem Fall auch
$p$, $q$, $r$ keinen gemeinsamen Teiler haben können, muss
auch $2$ durch $D$ teilbar sein. Dies ist nur möglich, wenn
$D$ entweder $=\pm1$ oder $=\pm2$ ist. Für alle Werte von $D$,
die größer als $2$ sind, ist daher stets $n=0$, wenn $\Delta$ zu
$D$ teilerfremd ist.

III.\ Ist dies geschehen, so muss offenbar jede eigentliche Darstellung
der Form $\varphi$ durch $f$ zu einem der übrig gebliebenen Werte gehören,
und zwar zu genau einem. Man muss diese Werte also der Reihe nach durchgehen
und die zu jedem gehörenden Darstellungen bestimmen. Um die Darstellungen
zu finden, die zu einem \textit{gegebenen} Wert $(B,B')$ gehören, bestimme
man zuerst die ternäre Form $g=\left(\begin{smallmatrix}a,\,a',\,a''\\b,\,b',\,b''\end{smallmatrix}\right)$ mit der Determinante $=\Delta$, in der
$a=p$, $b''=q$, $a'=r$, $ab-b'b''=B$, $a'b'-bb''=B'$ gilt. Die Werte von $a''$,
$b$, $b'$ erhält man daraus mithilfe der Gleichungen in II von
Art.\ 276. Aus diesen sieht man leicht, dass in dem Fall, wo $\Delta$,
$D$ teilerfremd sind, $b,b'\,a''$ notwendig ganze Zahlen sind (denn diese
drei Zahlen ergeben sowohl mit $D$ als auch mit $\Delta$ multipliziert
ganze Zahlen). Ist nun entweder einer der Koeffizienten $b,b',b''$ gebrochen,
oder sind die Formen $f$, $g$ nicht äquivalent, so kann es keine
Darstellungen der Form $\varphi$ durch $f$ geben, die zu $(B,B')$ gehören.
Sind dagegen $b,b',a''$ ganze Zahlen und die Formen $f$, $g$ äquivalent,
so liefert jede Transformation der ersten in die zweite, etwa
\[
\begin{array}{ccc}
\alpha,&\beta,&\gamma\\
\alpha',&\beta',&\gamma'\\
\alpha'',&\beta'',&\gamma''
\end{array}
\]
eine solche Darstellung, nämlich
\[
x=\alpha t+\beta u,\quad x'=\alpha't+\beta'u,\quad x''=\alpha''t+\beta''u
\]
Offenbar kann es keine Darstellung dieser Art geben, die sich nicht aus
einer Transformation ableiten ließe. So ist nun der Teil des zweiten Problems,
der die \textit{eigentlichen} Darstellungen sucht, auf das dritte Problem
zurückgeführt.
\SourceLeaf{333}

IV.\ Verschiedene Transformationen der Form $f$ in $g$ liefern
stets verschiedene Darstellungen, außer in dem einzigen Fall, wo der Wert
$(B,B')$ sich selbst entgegengesetzt ist. In diesem Fall liefern immer
zwei Transformationen eine einzige Darstellung. Nehmen wir nämlich an,
$f$ gehe auch in $g$ über durch die Substitution
\[
\begin{array}{ccc}
\alpha,&\beta,&\delta\\
\alpha',&\beta',&\delta'\\
\alpha'',&\beta'',&\delta''
\end{array}
\]
(die dieselbe Darstellung wie die vorige Transformation
liefert). Bezeichnen wir mit $k$, $\mathfrak{k}$, $\zeta$, $\eta$ dieselben
Zahlen wie in II des vorigen Artikels, so gilt
\[
B=k\mathfrak{k}B+\eta\mathfrak{k}D,\quad B'=k\mathfrak{k}B'+\zeta\mathfrak{k}D
\]
Nimmt man also an, dass $k$, $\mathfrak{k}$ entweder beide $=+1$ oder
beide $=-1$ sind, so ist (da wir den Fall $D=0$ ausgeschlossen haben)
$\zeta=0$, $\eta=0$. Daraus folgt leicht $\delta=\gamma$, $\delta'=\gamma'$, $\delta''=\gamma''$.
Die beiden Transformationen können also nur dann verschieden sein, wenn
von den Zahlen $k$, $\mathfrak{k}$ die eine $+1$, die andere $-1$
ist. Dann gilt $B\equiv-B$, $B'\equiv-B'$, $(\mathrm{mod.}\,D)$, oder der Wert $(B,B')$ ist
sich selbst entgegengesetzt.

V.\ Aus den oben (Art.\ 271) angegebenen Kriterien für definite und
indefinite Formen folgt leicht: Ist $\Delta$ positiv, $D$ negativ
und $\varphi$ eine negative Form, so ist $g$ eine negative definite
Form. Ist dagegen $\Delta$ positiv und entweder $D$ positiv oder
$D$ negativ und $\varphi$ eine positive Form, so ist $g$ eine
indefinite Form. Da nun $f$, $g$ sicher nur dann äquivalent
sein können, wenn sie in dieser Eigenschaft übereinstimmen, können offenbar
binäre Formen mit positiver Determinante sowie positive binäre Formen
nicht durch eine negative ternäre Form eigentlich dargestellt werden;
ebenso wenig negative binäre Formen durch eine indefinite ternäre Form
mit positiver Determinante. Eine ternäre Form der ersten bzw.\ zweiten
Art stellt also nur binäre Formen der zweiten bzw.\ ersten Art dar.
Ebenso folgt, dass eine definite ternäre Form mit negativer Determinante
(\textit{d.h.}, eine positive Form) nur positive binäre Formen darstellt,
eine indefinite dagegen nur negative binäre Formen und Formen mit
positiver Determinante.

\Needspace{8\baselineskip}
\Article{284}
Da die \textit{uneigentlichen} Darstellungen der binären Form $\varphi$
mit der Determinante $D$ durch die ternäre Form $f$, der die
Form $F$ adjungiert ist, diejenigen sind, aus denen uneigentliche
Darstellungen
\SourceLeaf{334}
der Zahl $D$ durch die Form $F$ folgen, kann $\varphi$ offenbar
nur dann durch $f$ uneigentlich dargestellt werden, wenn $D$
Quadratfaktoren enthält. Nehmen wir an, alle Quadrate, die $D$ teilen
(außer $1$), seien $ee$, $e'e'$, $e''e''$ usw.\ (ihre Anzahl
ist endlich, da wir voraussetzen, dass nicht $D=0$ ist). Dann liefert
jede uneigentliche Darstellung der Form $\varphi$ durch $f$ eine
Darstellung der Zahl $D$ durch $F$, in der die Werte der
Unbestimmten eine der Zahlen $e$, $e'$, $e''$ usw.\ als größten
gemeinsamen Teiler haben. In diesem Sinn sagen wir der Kürze halber,
dass jede uneigentliche Darstellung der Form $\varphi$ zum Quadratteiler
$ee$ oder $e'e'$ oder $e''e''$ usw.\ \textit{gehört}. Alle Darstellungen
der Form $\varphi$, die zu demselben \textit{gegebenen} Quadratteiler
$ee$ gehören (dessen Wurzel $e$ wir positiv nehmen), findet man
nun nach den folgenden Regeln. Aus ihrem synthetischen Beweis, den wir
hier der Kürze halber vorziehen, lässt sich leicht die Analysis wiederherstellen,
durch die sie gefunden wurden.

\textit{Erstens} bestimme man alle binären Formen mit der Determinante
$\frac{D}{ee}$, die in die Form $\varphi$ übergehen, und zwar durch eine eigentliche
Substitution der Form $T=\kappa t+\lambda u$, $U=\mu u$. Dabei bezeichnen $T$,
$U$ die Unbestimmten einer solchen Form; $t$, $u$ die
Unbestimmten der Form $\varphi$; $\kappa$, $\mu$ positive ganze Zahlen
(deren Produkt also $=e$ ist); $\lambda$ eine positive ganze Zahl
kleiner als $\mu$ (oder auch null). Diese Formen findet man mit den
zugehörigen Transformationen folgendermaßen:

Man setze $\kappa$ der Reihe nach jedem positiven Teiler von $e$
gleich (einschließlich $1$ und $e$) und $\mu=\frac{e}{\kappa}$. Für jedes
Paar bestimmter Werte von $\kappa$, $\mu$ gebe man $\lambda$ alle
ganzzahligen Werte von $0$ bis $\mu-1$. So erhält man sicher alle
Transformationen. Die Form, die durch eine beliebige Substitution
$T=\kappa t+\lambda u$, $U=\mu u$ in $\varphi$ übergeht, findet man nun, indem man die
Form bestimmt, in die $\varphi$ durch $t=\frac{1}{\kappa}T-\frac{\lambda}{e}U$, $u=\frac{1}{\mu}U$ übergeht. So
erhält man die zu den einzelnen Transformationen gehörenden Formen;
von allen diesen behalte man aber nur diejenigen, in denen alle drei
Koeffizienten ganze Zahlen sind\HistoricalNote{*)}{Dürften wir dieses Problem hier ausführlicher behandeln, so könnten wir die Lösung erheblich verkürzen. Es ist sofort klar, dass man für $\kappa$ nur solche Teiler von $e$ nehmen muss, deren Quadrat den ersten Koeffizienten der Form $\varphi$ teilt. Dieses Problem, aus dem sich auch einfachere Lösungen der Probleme in Art.\ 213, 214 ableiten lassen, behalten wir uns für eine andere geeignete Gelegenheit vor.}.

\textit{Zweitens} sei $\Phi$ eine dieser Formen, die in $\varphi$ durch
die Substitution $T=\kappa t+\lambda u$, $U=\mu u$ übergeht. Man suche alle \textit{eigentlichen}
\SourceLeaf{335}
Darstellungen der Form $\Phi$ durch $f$ (falls es solche gibt)
und bezeichne sie allgemein durch
\[
x=\mathfrak{A}T+\mathfrak{B}U,\quad x'=\mathfrak{A}'T+\mathfrak{B}'U,\quad x''=\mathfrak{A}''T+\mathfrak{B}''U\quad\ldots\quad(\mathfrak{R})
\]
Schließlich leite man aus jeder $(\mathfrak{R})$ die Darstellung
\[
x=\alpha t+\beta u,\quad x'=\alpha't+\beta'u,\quad x''=\alpha''t+\beta''u\quad\ldots\quad(\rho)
\]
mithilfe der Gleichungen
\[
\begin{array}{llll}
\alpha=\kappa\mathfrak{A},&\alpha'=\kappa\mathfrak{A}',&\alpha''=\kappa\mathfrak{A}''&\ldots\quad(R)\\
\beta=\lambda\mathfrak{A}+\mu\mathfrak{B},&\beta'=\lambda\mathfrak{A}'+\mu\mathfrak{B}',&\beta''=\lambda\mathfrak{A}''+\mu\mathfrak{B}''&
\end{array}
\]
ab. Man behandle die übrigen nach der ersten Regel gefundenen Formen
(falls es weitere gibt) genauso wie die Form $\Phi$, sodass aus jeder
ihrer eigentlichen Darstellungen weitere Darstellungen abgeleitet werden.
Ich behaupte, dass man so sämtliche Darstellungen der Form $\varphi$
erhält, die zum Teiler $ee$ gehören, und zwar jede genau einmal.

\textit{Beweis.} I.\ Dass die ternäre Form $f$ durch jede Substitution
$(\rho)$ tatsächlich in $\varphi$ übergeht, ist so offensichtlich, dass
keine weitere Erklärung nötig ist. Dass aber jede Darstellung $(\rho)$
uneigentlich ist und zum Teiler $ee$ gehört, folgt daraus, dass die
Zahlen $\alpha'\beta''-\alpha''\beta'$, $\alpha''\beta-\alpha\beta''$, $\alpha\beta'-\alpha'\beta$ jeweils $=e(\mathfrak{A}'\mathfrak{B}''-\mathfrak{A}''\mathfrak{B}')$, $e(\mathfrak{A}''\mathfrak{B}-\mathfrak{A}\mathfrak{B}'')$, $e(\mathfrak{A}\mathfrak{B}'-\mathfrak{A}'\mathfrak{B})$
werden. Ihr größter gemeinsamer Teiler ist also offenbar $e$ (weil
$(\mathfrak{R})$ eine eigentliche Darstellung ist).

II.\ Wir werden zeigen, dass sich aus jeder gegebenen Darstellung
$(\rho)$ der Form $\varphi$ eine eigentliche Darstellung einer Form mit
der Determinante $\frac{D}{ee}$ finden lässt, die unter den nach der ersten
Regel gefundenen Formen enthalten ist. Mit anderen Worten: Aus den
gegebenen Werten von $\alpha,\allowbreak \alpha',\allowbreak \alpha'',\allowbreak \beta,\allowbreak \beta',\allowbreak \beta''$ lassen sich ganzzahlige Werte von $\kappa$,
$\lambda$, $\mu$, die den vorgeschriebenen Bedingungen genügen, und
Werte von $\mathfrak{A}$, $\mathfrak{A}'$, $\mathfrak{A}''$, $\mathfrak{B}$, $\mathfrak{B}'$, $\mathfrak{B}''$,
die den Gleichungen $(R)$ genügen, ableiten, und zwar auf genau eine
Art. Zunächst folgt sofort aus den ersten drei Gleichungen in $(R)$,
dass man für $\kappa$ den positiv genommenen größten gemeinsamen Teiler
von $\alpha$, $\alpha'$, $\alpha''$ nehmen muss (denn da $\mathfrak{A}'\mathfrak{B}''-\mathfrak{A}''\mathfrak{B}'$, $\mathfrak{A}''\mathfrak{B}-\mathfrak{A}\mathfrak{B}''$,
$\mathfrak{A}\mathfrak{B}'-\mathfrak{A}'\mathfrak{B}$ keinen gemeinsamen Teiler haben dürfen, können auch $\mathfrak{A}$,
$\mathfrak{A}'$, $\mathfrak{A}''$ keinen haben). Damit sind auch $\mathfrak{A}$, $\mathfrak{A}'$,
$\mathfrak{A}''$ bestimmt, ebenso $\mu=\frac{e}{\kappa}$ (das notwendig ganzzahlig ist, wie
man leicht sieht). Nehmen wir an, drei ganze Zahlen $\mathfrak{a}$, $\mathfrak{a}'$,
$\mathfrak{a}''$ seien so gewählt, dass $\mathfrak{a}\mathfrak{A}+\mathfrak{a}'\mathfrak{A}'+\mathfrak{a}''\mathfrak{A}''=1$ gilt; und schreiben wir der
Kürze halber $k$ für $\mathfrak{a}\mathfrak{B}+\mathfrak{a}'\mathfrak{B}'+\mathfrak{a}''\mathfrak{B}''$.
\SourceLeaf{336}
Dann folgt aus den letzten drei Gleichungen $(R)$, dass $\mathfrak{a}\beta+\mathfrak{a}'\beta'+\mathfrak{a}''\beta''=\lambda+\mu k$
gelten muss. Daraus sieht man sofort, dass es für $\lambda$ genau einen
Wert zwischen den Grenzen $0$ und $\mu-1$ gibt. Da damit auch
$\mathfrak{B}$, $\mathfrak{B}'$, $\mathfrak{B}''$ bestimmte Werte erhalten, bleibt nur zu
zeigen, dass diese stets ganze Zahlen sind. Es gilt aber
\[
\begin{aligned}
\mathfrak{B}=\frac{1}{\mu}(\beta-\lambda\mathfrak{A})
&=\frac{1}{\mu}(\beta(1-\mathfrak{a}\mathfrak{A})-\mathfrak{A}(\mathfrak{a}'\beta'+\mathfrak{a}''\beta''))+\mathfrak{A}k\\
&=\frac{1}{\mu}(\mathfrak{a}''(\mathfrak{A}''\beta-\mathfrak{A}\beta'')-\mathfrak{a}'(\mathfrak{A}\beta'-\mathfrak{A}'\beta))+\mathfrak{A}k\\
&=\frac{1}{e}(\mathfrak{a}''(\alpha''\beta-\alpha\beta'')-\mathfrak{a}'(\alpha\beta'-\alpha'\beta))+\mathfrak{A}k
\end{aligned}
\]
Dieser Wert ist daher offenbar ganzzahlig. Ebenso bestätigt man leicht, dass
auch $\mathfrak{B}'$, $\mathfrak{B}''$ ganzzahlige Werte erhalten. --- Aus diesen
Überlegungen folgt, dass keine uneigentliche Darstellung der Form
$\varphi$ durch $f$, die zum Teiler $ee$ gehört, existieren kann,
die man mit der angegebenen Methode entweder gar nicht oder mehrmals erhält.

Behandelt man nun die übrigen Quadratteiler von $D$ genauso und
bestimmt die zu jedem gehörenden Darstellungen, so erhält man sämtliche
uneigentlichen Darstellungen der Form $\varphi$ durch $f$.

Aus dieser Lösung folgt ferner leicht, dass der am Ende des vorigen
Artikels für eigentliche Darstellungen angegebene Satz auch für uneigentliche
gilt: Allgemein lässt sich keine positive binäre Form mit negativer
Determinante durch eine negative ternäre Form darstellen usw.\ Ist
nämlich $\varphi$ eine solche binäre Form, die wegen jenes Satzes nicht
durch $f$ eigentlich dargestellt werden kann, so können auch alle
Formen mit den Determinanten $\frac{D}{ee}$, $\frac{D}{e'e'}$ usw., die $\varphi$ enthalten,
nicht durch $f$ eigentlich dargestellt werden. Denn alle diese Formen
haben eine Determinante mit demselben Vorzeichen wie $\varphi$. Sind
diese Determinanten negativ, so sind sie alle positive oder alle negative
Formen, je nachdem $\varphi$ zur ersten oder zur zweiten Art gehört.

